\documentclass[10pt,a4paper]{amsart}
\usepackage{amsmath,amssymb,amsthm,graphicx}
% Retained preamble for the Stacks Project chapter snapshot in this directory.
% The chapter's own preamble.tex is a copy of the shared Stacks Project preamble; the block
% below is sliced MECHANICALLY, line for line, from it. Nothing was added, removed or
% redefined; the wrapper class, the xr-hyper/hyperref wiring and the externaldocument list
% are omitted because the assembler supplies the document class.
%
% lines taken from preamble.tex: 3, 5, 9, 11, 12, 156, 157, 158, ... (51 lines total)
%
% package coverage restored from the chapter's own preamble (2026-10-13): xy, xy-2cell

\usepackage[all]{xy}
\xyoption{2cell}
\UseAllTwocells

\newenvironment{reference}{\comment}{\endcomment}
\newenvironment{slogan}{\comment}{\endcomment}
\newenvironment{history}{\comment}{\endcomment}
\theoremstyle{plain}
\newtheorem{theorem}[subsection]{Theorem}
\newtheorem{proposition}[subsection]{Proposition}
\newtheorem{lemma}[subsection]{Lemma}
\theoremstyle{definition}
\newtheorem{definition}[subsection]{Definition}
\newtheorem{example}[subsection]{Example}
\newtheorem{exercise}[subsection]{Exercise}
\newtheorem{situation}[subsection]{Situation}
\theoremstyle{remark}
\newtheorem{remark}[subsection]{Remark}
\newtheorem{remarks}[subsection]{Remarks}
\numberwithin{equation}{subsection}
\def\lim{\mathop{\mathrm{lim}}\nolimits}
\def\colim{\mathop{\mathrm{colim}}\nolimits}
\def\Spec{\mathop{\mathrm{Spec}}}
\def\Hom{\mathop{\mathrm{Hom}}\nolimits}
\def\Ext{\mathop{\mathrm{Ext}}\nolimits}
\def\SheafHom{\mathop{\mathcal{H}\!\mathit{om}}\nolimits}
\def\SheafExt{\mathop{\mathcal{E}\!\mathit{xt}}\nolimits}
\def\Sch{\mathit{Sch}}
\def\Mor{\mathop{\mathrm{Mor}}\nolimits}
\def\Ob{\mathop{\mathrm{Ob}}\nolimits}
\def\Sh{\mathop{\mathit{Sh}}\nolimits}
\def\NL{\mathop{N\!L}\nolimits}
\def\CH{\mathop{\mathrm{CH}}\nolimits}
\def\proetale{{pro\text{-}\acute{e}tale}}
\def\etale{{\acute{e}tale}}
\def\QCoh{\mathit{QCoh}}
\def\Ker{\mathop{\mathrm{Ker}}}
\def\Im{\mathop{\mathrm{Im}}}
\def\Coker{\mathop{\mathrm{Coker}}}
\def\Coim{\mathop{\mathrm{Coim}}}
\DeclareMathSymbol{\boxtimes}{\mathbin}{AMSa}{"02}
\def\QCohstack{\mathcal{QC}\!\mathit{oh}}
\def\Cohstack{\mathcal{C}\!\mathit{oh}}
\def\Spacesstack{\mathcal{S}\!\mathit{paces}}
\def\Quotfunctor{\mathrm{Quot}}
\def\Hilbfunctor{\mathrm{Hilb}}
\def\Curvesstack{\mathcal{C}\!\mathit{urves}}
\def\Polarizedstack{\mathcal{P}\!\mathit{olarized}}
\def\Complexesstack{\mathcal{C}\!\mathit{omplexes}}
\def\Pic{\mathop{\mathrm{Pic}}\nolimits}
\def\Picardstack{\mathcal{P}\!\mathit{ic}}
\def\Picardfunctor{\mathrm{Pic}}
\def\Deformationcategory{\mathcal{D}\!\mathit{ef}}
\usepackage{cleveref}
\crefname{originallemma}{original lemma}{original lemmas}
\Crefname{originallemma}{Original lemma}{Original lemmas}
\makeatletter
\newlabel{categories-lemma-functorial-limit}{{002L}{source}{}{sourcepointer-1}{}}
\newlabel{categories-lemma-functorial-limit@cref}{{[originallemma][0][]002L}{[1][0][]source}{}{}{}}
\makeatother
\title{Stacks, Sheaves on Algebraic Stacks: lemma-1-morphisms-presheaves}
\author{The Stacks Project Authors}
\date{stacks-project:stacks-sheaves.tex:local-snapshot-20261012; GNU FDL 1.2 or later}
\begin{document}
\maketitle
\noindent\textit{Editorial note.} Statement span is lines 369-378 of sources/stacks-algebraic-stacks-sheaves/pinned-original-source.tex; proof span is lines 380-437, the complete author proof (56 body lines).\par\medskip
\noindent\textit{Editorial note.} Navigation stubs: 1 cited-but-not-included labels are supplied as non-load-bearing pointers with their pinned tags where the snapshot tag index has them.\par\medskip
\noindent\textit{Editorial note.} extra\_packages is empty on purpose: the retained preamble already loads xy with options and a second declaration clashes in the assembler wrapper.\par\medskip
\noindent\textit{Editorial note.} Text is retained verbatim, including the author's own typos and placeholders.\par\medskip
\noindent\textit{Editorial note.} Queue index 9 of queue\_algebraic-stacks-sheaves.json (pb=2136, stmt 369-378, proof 380-437).\par\medskip
\noindent\textit{Allowed standard prerequisites (editorial).}\begin{itemize}
\item Explicit prerequisite boundary. The cited inputs are prerequisites; the argument in the retained span is not replaced by any cited result.
\end{itemize}\medskip
\ifdefined\hypertarget\hypertarget{sourcepointer-1}{}\fi
\noindent\textit{Source reference.} Original lemma 002L: sources/stacks-algebraic-stacks-sheaves/pinned-original-source.tex, source lines 437--437. Cited by the retained statement or proof but not part of the collected spans; supplied only as a navigation pointer. The argument does not depend on its content.\par\medskip
\begin{lemma}
\label{lemma-1-morphisms-presheaves}
Let $f : \mathcal{X} \to \mathcal{Y}$ and $g : \mathcal{Y} \to \mathcal{Z}$
be $1$-morphisms of categories fibred in groupoids over
$(\Sch/S)_{fppf}$. Then $(g \circ f)^p = f^p \circ g^p$ and
there is a canonical isomorphism
${}_p(g \circ f) \to {}_pg \circ {}_pf$
compatible with adjointness of $(f^p, {}_pf)$, $(g^p, {}_pg)$, and
$((g \circ f)^p, {}_p(g \circ f))$.
\end{lemma}

\begin{proof}
Let $\mathcal{H}$ be a presheaf on $\mathcal{Z}$. Then
$(g \circ f)^p\mathcal{H} = f^p (g^p\mathcal{H})$ is given
by the equalities
$$
(g \circ f)^p\mathcal{H}(x) = \mathcal{H}((g \circ f)(x))
= \mathcal{H}(g(f(x))) = f^p (g^p\mathcal{H})(x).
$$
We omit the verification that this is compatible with restriction maps.

\medskip\noindent
Next, we define the transformation ${}_p(g \circ f) \to {}_pg \circ {}_pf$.
Let $\mathcal{F}$ be a presheaf on $\mathcal{X}$.
If $z$ is an object of $\mathcal{Z}$ then we get a
category $\mathcal{J}$ of quadruples
$(x, f(x) \to y, y, g(y) \to z)$ and a category $\mathcal{I}$
of pairs $(x, g(f(x)) \to z)$. There is a canonical functor
$\mathcal{J} \to \mathcal{I}$ sending the object
$(x, \alpha : f(x) \to y, y, \beta : g(y) \to z)$ to
$(x, \beta \circ f(\alpha) : g(f(x)) \to z)$. This gives the arrow in
\begin{align*}
({}_p(g \circ f)\mathcal{F})(z) & =
\lim_{g(f(x)) \to z} \mathcal{F}(x) \\
& = \lim_\mathcal{I} \mathcal{F} \\
& \to \lim_\mathcal{J} \mathcal{F} \\
& = \lim_{g(y) \to z}
\Big(\lim_{f(x) \to y} \mathcal{F}(x)\Big) \\
& =
({}_pg \circ {}_pf\mathcal{F})(x)
\end{align*}
by
Categories, Lemma \ref{categories-lemma-functorial-limit}.
We omit the verification that this is compatible with restriction maps.
An alternative to this direct construction is to define
${}_p(g \circ f) \cong {}_pg \circ {}_pf$
as the unique map compatible with the adjointness properties. This also
has the advantage that one does not need to prove the compatibility.

\medskip\noindent
Compatibility with adjointness of $(f^p, {}_pf)$, $(g^p, {}_pg)$, and
$((g \circ f)^p, {}_p(g \circ f))$ means that given presheaves
$\mathcal{H}$ and $\mathcal{F}$ as above we have a commutative diagram
$$
\xymatrix{
\Mor_{\textit{PSh}(\mathcal{X})}(f^pg^p\mathcal{H}, \mathcal{F})
\ar@{=}[r] \ar@{=}[d] &
\Mor_{\textit{PSh}(\mathcal{Y})}(g^p\mathcal{H}, {}_pf\mathcal{F})
\ar@{=}[r] &
\Mor_{\textit{PSh}(\mathcal{Y})}(\mathcal{H}, {}_pg{}_pf\mathcal{F})
\\
\Mor_{\textit{PSh}(\mathcal{X})}((g \circ f)^p\mathcal{G}, \mathcal{F})
\ar@{=}[rr] & &
\Mor_{\textit{PSh}(\mathcal{Y})}(\mathcal{G}, {}_p(g \circ f)\mathcal{F})
\ar[u]
}
$$
Proof omitted.
\end{proof}

\end{document}
